% !TeX root = ../初等数学笔记.tex
% 由主文档通过 \input 载入；请编译主文档。
\setcounter{chapter}{0} % 保留原书章号。
\chapter{符号、公式与方法索引}
\label{app:em-notation}
本附录用于定位工具。公式不能脱离条件单独使用，具体证明和例题应返回对应章节阅读。

\section{符号与中英文术语}
$\mathbb N$ 包含零，$\mathbb N_{>0}$ 表示正整数；$\mathbb Z,\mathbb Q,\mathbb R,\mathbb C$ 依次为整数、有理数、实数和复数集。$A\subseteq B$ 表示集合包含，$x\in A$ 表示元素关系；$|A|$ 在有限集合语境中是基数（cardinality），$|x|$ 对实数是绝对值，对复数是模。$\gcd$ 表示最大公因数，$\operatorname{lcm}$ 表示最小公倍数。

$f^{-1}$ 表示单射函数在值域上的反函数，不表示倒数。$a_n$ 是数列第 $n$ 项，$S_n$ 是前 $n$ 项和，$\binom nk$ 表示组合数。$\boldsymbol u\cdot\boldsymbol v$ 为数量积，$\overline z$ 为复共轭；统计中的 $\overline x$ 则表示样本均值，含义由所定义对象确定。$\mathbb P,\mathbb E,\operatorname{Var}$ 分别表示概率、期望和方差，随机变量用大写字母，观测值通常用小写。

\section{常用公式及其条件}
代数求解可查二次方程公式 \eqref{eq:em-quadratic}、因式定理 \eqref{eq:em-factor} 和 Bézout 等式 \eqref{eq:em-bezout}；求和可查等比和 \eqref{eq:em-geometric-sum} 与二项式定理 \eqref{eq:em-binomial}。不等式的变量范围和取等条件见第\ref{chap:em-inequality-methods}章，特别是式 \eqref{eq:em-cauchy}。

三角推导的起点为平方关系 \eqref{eq:em-trig-pythagorean} 与和差角公式 \eqref{eq:em-cos-add}、\eqref{eq:em-sin-add}，解三角形用式 \eqref{eq:em-cosine-rule}。坐标距离用式 \eqref{eq:em-point-line}，三类圆锥曲线的参数约定分别见式 \eqref{eq:em-ellipse}、\eqref{eq:em-hyperbola}、\eqref{eq:em-parabola}。概率分来源计算用式 \eqref{eq:em-total-probability}，使用前须检查分割完备性。

\section{问题结构与方法检索}
多项式零点先看因式分解和余式，分式符号先找全部零点与禁用点，含绝对值表达式先用距离或分段。定和定积约束优先检验均值不等式和配方，重复累积先识别等差、等比或可平移的递推。几何中的等长与比例联系全等、相似与点幂，位置和方向可转化为向量或坐标。

有限选择先明确顺序、重复及对象等价关系，再用加法、乘法或容斥；存在性可尝试抽屉原理，不可达性可寻找不变量。概率先建样本空间和权重，统计先查抽样口径与单位。每个方法都应以满足其适用条件为入口，关键词相似本身不构成使用依据。

